\begin{lemma}[Part IV, Algorithm (C) and invariant (24)]
Fix the following Case 2 data.  Let \(v_n\in V(T)\) be the terminal
supporting vertex, and let
\[
 L_1=\noderef{ExtensionBaseSupportClass}(v_n)
      =\{1_1<1_2<\cdots<1_{p+1}\}
\]
be the exact increasing enumeration of the full terminal support class.
Assume that its last member is the terminal slot
\(1_{p+1}=n\), that \(n+1\) is the number of certified rounds, and
that every displayed slot has PE kind, exact supporting vertex \(v_n\),
exact certified fields \(\gamma_{1_a},\delta_{1_a}\), and exact stored
set \(S_{1_a}=\{\gamma_{1_a},\delta_{1_a}\}\).  Retain
explicitly the support-color conclusions furnished in the extension-base
argument by \noderef{ExtensionBaseSupportClassColors}:
\[
 \gamma_{1_{a+1}}=\delta_{1_a}\quad(1\leq a\leq p),
 \qquad
 \gamma_{1_1},\delta_{1_1},\ldots,\delta_{1_{p+1}}
 \text{ form an injectively indexed color chain}.             \tag{1}
\]
Let \(T_j\) be the certified tree stored at an exact slot \(j\), and
assume \(T_j\) is an ordered initial segment, in the sense of
\noderef{TreeSequenceInitialSegment}, of \(T_t\) for every
\(t\in L_1\setminus\{n\}\).

Let \(\mu_0\) be proper and leave exactly \(e_0\) uncolored.  Assume
the following explicit consequences of the construction in Claim 4.6:
\[
 \gamma_n\in\overline\mu_0(u),                               \tag{2}
\]
and, for every exact \(t\in L_1\setminus\{n\}\), the pair
\(\gamma_t,\delta_t\) is \(T_t\)-strongly interchangeable under
\(\mu_0\).  Also retain the present-color fact implicit in that
construction: every preterminal \(\gamma_t\) is present at \(u\) under
\(\mu_0\).  Finally assume the precise form of (25), namely
\[
 \overline\mu_0(v_n)\cap
   \bigcup_{t\in L_1\setminus\{n\}}S_t=\varnothing,           \tag{3}
\]
where \(S_t\) is the exact stored-color set at slot \(t\).

Then there is an \noderef{ExtensionBaseAlgorithmCTrace} on
\(L_1^*=L_1\setminus\{n\}\) satisfying
\noderef{ExtensionBaseAlgorithmCInvariant}.  Thus Algorithm (C) processes
the exact slots of \(L_1^*\) in decreasing order, starting at \(\mu_0\),
and invariant (24) holds before and after every rooted component switch.
This includes the zero-step case \(p=0\).
\end{lemma}
\begin{proof}
Write the increasing enumeration as
\(i_a=1_{a+1}\) for \(0\leq a\leq p\), so \(i_p=n\), and put
\[
 A=L_1\setminus\{n\}=\{i_0,\ldots,i_{p-1}\}.                \tag{4}
\]
The equality in (4) is an equality of exact finite slot identities.  For
the forward inclusion, every \(i_a\) belongs to the exact enumerated
support class, and strict increase gives \(i_a<i_p=n\) when \(a<p\).
For the reverse inclusion, exactness of the enumeration writes a member
of \(L_1\) uniquely as \(i_q\); if it is not \(n=i_p\), strict increase
forces \(q<p\).  Here and below we use
\noderef{ExtensionBaseSupportClass} with the exact range hypothesis, not
merely equality of numerical slot values.

If \(p=0\), (4) says that \(A\) is empty.  Take the zero-step
\noderef{ExtensionBaseAlgorithmCTrace}, whose only coloring is
\(\mu_0\).  Its terminal slot is \(i_0=n\), which belongs to \(L_1\),
and the terminal-round and initial-color fields are hypotheses.  The
properness and unique-uncolored clauses of
\noderef{ExtensionBaseAlgorithmCInvariant} are the two initial coloring
hypotheses.  Every remaining clause quantifies over a step or a member of
the empty set \(A\), so it is vacuous.  This proves the assertion when
\(p=0\).  Henceforth assume \(p>0\).

For \(0\leq a<p\), define
\[
 r_a=p-1-a,\qquad t_a=i_{r_a}.                                \tag{5}
\]
The list \(t_0,\ldots,t_{p-1}\) has exact image \(A\) by (4).  If
\(a<c<p\), then \(r_c<r_a\), and strict increase of the original
enumeration gives \(t_c<t_a\).  Define recursively
\[
 \mu^0=\mu_0,\qquad
 K_a=\noderef{TwoColorComponent}
   (G,\mu^a,\gamma_{r_a},\delta_{r_a},\mathsf{inl}(u)),
 \qquad
 \mu^{a+1}=\noderef{TwoColorComponentSwap}
   (G,\mu^a,\gamma_{r_a},\delta_{r_a},\mathsf{inl}(u)).       \tag{6}
\]
At step \(a\), use the exact slot \(t_a\), the exact colors
\(\gamma_{r_a},\delta_{r_a}\), and the coloring \(\mu^a\).  Equations
(4)--(6), together with the assumed support, PE-kind, gamma, and delta
field equalities, verify every field of an
\noderef{ExtensionBaseAlgorithmCTrace}; call it \(\mathcal C\).

For the color bookkeeping, set
\[
 c_0=\gamma_0,\qquad c_{q+1}=\delta_q\quad(0\leq q\leq p).   \tag{7}
\]
The chain identities in (1) give \(\gamma_q=c_q\) for every
\(0\leq q\leq p\).  The injectivity in (1) says that
\(c_0,\ldots,c_{p+1}\) are pairwise distinct.  Consequently
\[
 \gamma_q=c_q\ne c_{q+1}=\delta_q,                           \tag{8}
\]
and, if \(q<r\), then \(\gamma_q=c_q\) differs from both colors
\(c_r,c_{r+1}\) switched at position \(r\).  This is precisely the
chain and distinctness conclusion stated by
\noderef{ExtensionBaseSupportClassColors}; no equality of two different
certified color fields is being inferred merely from their names.

Let
\[
 U=\bigcup_{t\in A}S_t.                                      \tag{9}
\]
For every \(q<p\), the exact stored-pair field at slot \(i_q\) puts
\(\gamma_q\) and \(\delta_q\) in \(U\).  Hypothesis (3), together
with the definitions \noderef{MissingColors} and
\noderef{PresentColors}, therefore gives an edge incident with \(v_n\)
of each of these two colors under \(\mu_0\).  We shall preserve this
stronger presentness assertion for every color in \(U\), rather than
using (3) only once.

We prove invariant (24) by induction along the decreasing list (5).
Immediately before processing position \(r=r_a\), maintain:
\[
\begin{array}{ll}
({\rm I})&\mu^a\text{ is proper and leaves exactly }e_0
          \text{ uncolored};\\
({\rm II})&\delta_r\in\overline{\mu^a}(u),\quad
          \gamma_q\text{ is present at }u\text{ for every }q\leq r;\\
({\rm III})&\text{every pair at an unprocessed slot is strongly
          interchangeable under }\mu^a;\\
({\rm IV})&\text{every color in }U\text{ is present at }v_n
          \text{ under }\mu^a.
\end{array}                                                    \tag{10}
\]
At \(a=0\), assertion (I) is assumed.  Since \(r_0=p-1\), the last
chain identity gives
\(\delta_{r_0}=\gamma_p=\gamma_n\), so (2) gives the missing-delta
part of (II); the explicitly assumed initial presentness of all
preterminal gammas gives its other part.  The strong-interchangeability
conclusion of Claim 4.6, stated here for every exact member of \(A\), is
(III).  The preceding paragraph is (IV).

Fix a stage \(a\), abbreviate
\[
 r=r_a,\quad t=t_a,\quad \alpha=\gamma_r,\quad
 \beta=\delta_r,\quad \varphi=\mu^a,\quad K=K_a.             \tag{11}
\]
We verify the per-step clauses of
\noderef{ExtensionBaseAlgorithmCInvariant} and then restore (10).

By (II), \(\beta\) is missing and \(\alpha\) is present at \(u\).
Unfolding \noderef{MissingColors} and \noderef{PresentColors}, no edge
incident with \(u\) has \(\varphi\)-color \(\beta\), while some edge
\(e_\alpha\) incident with \(u\) has color \(\alpha\).  This
\(\alpha\)-edge is unique.  Indeed, two distinct such edge identities
would share \(u\), hence would be adjacent by \noderef{AdjacentEdges},
contradicting properness as defined by
\noderef{IsProperEdgeColoring}.  The definition
\noderef{TwoColorIncidenceGraph} makes
\(\mathsf{inr}(e_\alpha)\) adjacent to \(\mathsf{inl}(u)\), so both
nodes belong to \(K\).  The preceding uniqueness and missingness show
that \(u\) has exactly one neighbor in this incidence graph.

Thus \(K\) contains a genuine edge node and contains a node with a
subsingleton neighbor set.  Distinctness (8) and
\noderef{TwoColorComponentSpanningPath}, applied to the proper coloring
in (I), show that \(K\) is exactly a path component in the sense of
\noderef{TwoColorPathComponent}, with \(u\) as an end.  This proves the
path clause, including the nondegeneracy that would fail if both colors
were missing at \(u\).

The current pair is strongly interchangeable by (III).  Suppose
\(x\in V(T_t)\), \(x\ne v_n\), and \(\mathsf{inl}(x)\in K\).  Since
\(x\) and \(u\) lie in the same connected component, the component
through \(x\) is exactly \(K\).  Expanding
\noderef{ETTStrongInterchangeable} at \(x\) supplies a
\noderef{TwoColorComponentCycle} whose support is all of \(K\), not
merely a cycle contained in it.  Its \noderef{TwoColorCycle} therefore
contains the endpoint \(u\).  The two directions of a simple closed
cycle at that occurrence of \(u\) give two distinct neighbors, contrary
to the singleton neighbor set just proved.  Hence
\[
 x\in V(T_t),\ \mathsf{inl}(x)\in K\quad\Longrightarrow\quad x=v_n.
                                                                    \tag{12}
\]
This is the exact ``at most \(v_n\)'' clause of (24).

By (9), (IV), and the exact stored-pair field at \(t\), both
\(\alpha\) and \(\beta\) are present at \(v_n\).  Thus there are
incident edge identities \(e_\alpha^n,e_\beta^n\) of those colors.
They are distinct by (8).  If \(\mathsf{inl}(v_n)\in K\), both are
neighbors of \(v_n\) in \(K\).  The degree bound
\noderef{TwoColorIncidenceDegreeAtMostTwo}, using properness in (I), says
that these two distinct neighbors exhaust the neighbor set.  In
particular it is not subsingleton, so \(v_n\) is not an end whenever it
lies on the path.  This is exactly the remaining path assertion in (24).

The same argument preserves (IV).  Fix \(c\in U\) and an incident
\(c\)-colored edge at \(v_n\).  If \(v_n\notin K\), no incident edge
at \(v_n\) can be changed by (6), since any changed edge-node would be
adjacent to and hence place \(v_n\) in \(K\).  If \(v_n\in K\), the
preceding paragraph shows that the swap exchanges the colors on exactly
\(e_\alpha^n,e_\beta^n\), so both current colors remain present.  Every
other color is fixed pointwise by
\noderef{TwoColorComponentSwap}.  Therefore every color in \(U\)
remains present at \(v_n\) under \(\mu^{a+1}\).

At \(u\), the unique incident \(\alpha\)-edge belongs to \(K\) and is
changed to \(\beta\); there was no incident \(\beta\)-edge.  Hence no
incident edge has color \(\alpha\) after the swap, while an incident edge
does have color \(\beta\).  Thus \(\alpha\) becomes missing and
\(\beta\) becomes present, by \noderef{Incident},
\noderef{MissingColors}, and \noderef{PresentColors}.  For every
\(q<r\), (7)--(8) show that \(\gamma_q=c_q\) differs from both switched
colors \(c_r,c_{r+1}\), so its existing incident witness at \(u\) is
fixed.  If a next position exists, it is \(r-1\), and (1) gives
\[
                  \delta_{r-1}=\gamma_r=\alpha.              \tag{13}
\]
The newly missing \(\alpha\), equation (13), and the surviving lower
gammas establish (II) for the next stage.

By (8), \(\alpha\ne\beta\), so
\noderef{TwoColorComponentSwapProper} applied to (I) proves that
\(\mu^{a+1}\) is proper.  The set of uncolored edge identities is
unchanged.  Indeed, expansion of
\noderef{TwoColorComponentSwap} gives, for every edge \(f\),
\[
 \mu^{a+1}(f)=\mathsf{none}\quad\Longleftrightarrow\quad
 \mu^a(f)=\mathsf{none}:                                    \tag{14}
\]
outside \(K\) the value is fixed, and inside \(K\) two
\(\mathsf{some}\)-values are exchanged while \(\mathsf{none}\) is
fixed.  Applying both clauses of \noderef{IsOneEdgeUncolored} to (14)
proves that precisely \(e_0\) remains uncolored.  This restores (I).

It remains to restore (III), which is the post-step assertion in (24).
Let \(s\in A\) be unprocessed after the current step, and let its exact
certified fields be \(\gamma_s,\delta_s\).  Exact enumeration gives a
unique \(q<r\) with \(s=i_q\); in particular strict increase gives
\(s<t\).  Comparing the certified field equalities with those at
position \(q\) identifies
\(\gamma_s=\gamma_q\) and \(\delta_s=\delta_q\).  The current and
surviving slots have the same exact support \(v_n\) and PE kind.  Their
pairs are strongly interchangeable under \(\mu^a\) by (III), and the
current rooted component is the path component already established.
Together with properness and unique-uncoloredness from (I), these are
all premises of the strict earlier-slot alternative in
\noderef{StrongInterchangeabilityKempePreserved}.  It follows that the
pair at \(s\) remains strongly interchangeable under \(\mu^{a+1}\).
The lemma's different-support disjointness premise is vacuous here:
both certified supporting vertices are exactly \(v_n\), so its required
alternative is supplied solely by the strict slot inequality \(s<t\).
Since \(s\) was arbitrary, (III) is restored for every unprocessed slot.

Finally, the assumed ordered initial-segment premise applied to the exact
current member \(t\in A\) gives \(T_j\preccurlyeq T_t\).  Together with
(I), the missing-delta statement (II), the path certificate, (12), the
nonendpoint argument at \(v_n\), and the survivor conclusion, this
verifies every per-step field of
\noderef{ExtensionBaseAlgorithmCInvariant}.  The induction gives
properness and exact-one-uncoloredness for the output coloring as well.
By the exact image equality in (4)--(5), no member of \(A\) is absent
from the complete trace, so the invariant's final absent-slot clause is
vacuous.  Hence \(\mathcal C\) satisfies invariant (24), as required.
\end{proof}
